\begin{center}
	\large\textbf{ABSTRAK}
\end{center}

\addcontentsline{toc}{chapter}{ABSTRAK}

\vspace{2ex}
\begingroup
% Menghilangkan padding
\setlength{\tabcolsep}{0pt}

\noindent
\begin{tabularx}{\textwidth}{l >{\centering}m{2em} X}
	Nama Mahasiswa    & : & \name{}         \\
	
	Judul Tugas Akhir & : & \tatitle{}      \\
	
	Pembimbing        & : & 1. \advisor{}   \\
	&   & 2. \coadvisor{} \\
\end{tabularx} 
\endgroup

% Ubah paragraf berikut dengan abstrak dari tugas akhir
Kemajuan teknologi \textit{Unmanned Aerial Vehicle} (UAV) atau drone memberikan peluang besar untuk mempercepat operasi \textit{Search and Rescue} (SAR) di area bencana yang sulit dijangkau. Penelitian ini bertujuan merancang dan mengimplementasikan sistem otonom pada drone Holybro S500 untuk mendeteksi manusia secara \textit{real-time} berbasis \textit{edge computing} menggunakan Raspberry Pi 5 dan kamera \textit{webcam}. Sistem mengintegrasikan \textit{Flight Controller} Pixhawk 6C dan kendali penerbangan via protokol MAVLink melalui pustaka MAVSDK-Python. Berdasarkan evaluasi, model YOLOv8n hasil \textit{fine-tuning} berformat ONNX dengan resolusi 320x320 piksel dipilih karena lebih konsisten mempertahankan \textit{bounding box} pada pengujian luar ruangan dibanding model \textit{pose estimation}. Perangkat lunak sistem dikembangkan menggunakan arsitektur \textit{multi-threading} asinkron, didukung distribusi daya stabil berbasis regulator UBEC 5A untuk mencegah kegagalan sistem akibat \textit{voltage sag}. Hasil pengujian terbang nyata memvalidasi kemampuan navigasi \textit{Scout Mission} secara kontinu melalui siklus \textit{state machine} yang meliputi tahap pemindaian, pendekatan, dokumentasi, dan perpindahan area tanpa perlu mendarat setiap kali mendokumentasikan target. Selain itu, sistem keselamatan berlapis mampu memitigasi anomali pembacaan sensor akibat \textit{barometric drift} saat terbang rendah. Penelitian ini mengembangkan purwarupa drone otonom yang mengintegrasikan persepsi visual, pengambilan keputusan otonom, dan transmisi pemantauan \textit{real-time} untuk mendukung misi pencarian korban.
% Ubah kata-kata berikut dengan kata kunci dari tugas akhir
Kata Kunci: Drone Otonom, Komputasi Edge, MAVSDK-Python, Raspberry Pi 5, Pencarian dan Penyelamatan, YOLOv8n